\hojaejercicio{Aplicaciones Teorema de Cauchy}
%ejercicio 1
\ejercicio{% Ejercicio 1
    Desarrolla en serie de potencias de la forma $\sum_{n=0}^{\infty}c_{n}z^{n}$ la función:
    \[
    f(z) = \frac{1}{1-2z}
    \]
    ¿Dónde converge la serie?}{
Se puede ver de dos maneras:

Tenemos que 
\[
f(z) = \frac{1}{1 - 2z} \in H(\mathbb{C} \setminus \{ 1/2 \})
\]
Es un resultado de una serie geométrica de razón $2z$.  
Entonces:
\[
f(z) = \frac{1}{1 - 2z} = \sum_{n=0}^{\infty} (2z)^n
\quad \text{que converge si } |z| < \frac{1}{2}.
\]

Otra forma: vemos que $f(z) \in H(\mathbb{C} \setminus \{1/2\})$, entonces, al ser $\mathbb{C} \setminus \{1/2\}$ un dominio y $f$ holomorfa en él, entonces
\[
\forall z_0 \in \mathbb{C}, \quad \exists (c_n)_{n=0}^{\infty} \subset \mathbb{C} \text{ tal que (por el teorema de Taylor):}
\]
\[
f(z) = \sum_{n=0}^{\infty} c_n (z - z_0)^n, \quad \forall z \in D(z_0, r)
\]

Por lo tanto, sea $z_0 = 0$:
\[
f(z) = \sum_{n=0}^{\infty} c_n z^n, \quad \forall z \in D(0,1/2)
\]
Además:
\[
c_n = \frac{f^{(n)}(0)}{n!} = 2^n 
\Rightarrow f(z) = \sum_{n=0}^{\infty} 2^n z^n, \quad z \in D(0,1/2)
\]

}
%ejercicio 2
\ejercicio{
    Usa el Teorema de Liouville para probar que si $f \in H(\mathbb{C})$ (es una función entera) y cumple las condiciones de periodicidad:
    \[
    f(z+2\pi) = f(z) \quad \text{y} \quad f(z+2\pi i) = f(z)
    \]
    entonces $f$ es constante.
}{
     Vamos a probar que $f \in H(\mathbb{C})$ está acotada:

Sea $A = \{ a + bi : a,b \in [0,2\pi] \}$. Entonces, como 
$f(z + 2\pi) = f(z)$ y $f(z + 2\pi i) = f(z)$, 
vamos a tener que:
\[
f(A) = f(\mathbb{C})
\]

Esto es debido a que, sea $z \in \mathbb{C}$, entonces 
$z = z_1 + z_2 i$, con $z_1,z_2 \in \mathbb{R}$.  
Por lo tanto, 
\[
z_1 = k_1 + 2\pi m_1, \quad z_2 = k_2 + 2\pi m_2
\]
con $k_1, k_2 \in [0,2\pi)$ y $m_1,m_2 \in \mathbb{Z}$ (por la división euclídea).

\[
\Rightarrow f(z) = f(z_1 + z_2 i) = f(k_1 + 2\pi m_1 + (k_2 + 2\pi m_2)i)
= f(k_1 + k_2 i) \in f(A)
\]

Por lo tanto, $f(\mathbb{C}) \subseteq f(A)$ y $f(A) \subseteq f(\mathbb{C}) \Rightarrow f(A)=f(\mathbb{C})$.

Como $f(A) = f(\mathbb{C})$ y $A$ es cerrado y acotado (compacto) y f continua $\Rightarrow f(A)$ es cerrado y acotado.  
$\Rightarrow f$ acotada en $\mathbb{C} \Rightarrow \exists M \in \mathbb{R}$ tal que $|f| \le M$ para todo $z \in \mathbb{C}$.  
Y por el teorema de Liouville, entonces $f$ es constante.
}
%ejercicio 3
\ejercicio{
    Evalúa las siguientes integrales:
    \begin{enumerate}
        \item $\displaystyle \int_{S(0;1)} e^{z} z^{-3} \, dz$
        \item $\displaystyle \int_{S(i;1)} \frac{e^{z}}{(z-i)\cos z} \, dz$
    \end{enumerate}
}{
    \[
\int_{S(0;i)} \frac{e^z}{z^3} \, dz = \frac{2\pi i}{2} \cdot e^0 = \pi i
\]

Sea $g(z) = e^z \in H(\mathbb{C})$ y $\overline{D}(0,1) \subset \mathbb{C}$.  
Por la fórmula integral de Cauchy para las derivadas:
\[
g''(0) = \frac{2!}{2\pi i} \int_{\partial D(0,1)} \frac{f(w)}{(w - 0)^3} \, dw
\]

2.
\[
\int_{S(i;1)} \frac{e^z}{(z - i) \cos z} \, dz = \frac{2\pi i \, e^i}{\cos i}
\]

Sea $g(z) = \frac{e^z}{\cos z} \in H(\mathbb{C} \setminus \{ \pi/2 + n\pi,\, n \in \mathbb{Z} \})$  
Como $\overline{D}(i;1) \subset \mathbb{C} \setminus \{ \pi/2 + n\pi,\, n \in \mathbb{Z} \}$,  
por la fórmula integral de Cauchy:
\[
2\pi i \, g(i) = \int_{\partial D(i;1)} \frac{g(w)}{(w - i)} \, dw
\]
}

%ejercicio 4
\ejercicio{
    Estudia si es entera la función definida por:
    \[
    f(z) = \begin{cases} 
    \frac{e^{iz}-1}{iz} & \text{si } z \neq 0 \\
    1 & \text{si } z = 0 
    \end{cases}
    \]
}{
    \[
f(z) =
\begin{cases}
\dfrac{e^{iz} - 1}{iz}, & \text{si } z \neq 0, \\[1.5em]
1, & \text{si } z = 0.
\end{cases}
\]

Vemos que $\dfrac{e^{iz} - 1}{iz} \in H(\mathbb{C} \setminus \{0\})$ puesto que el cociente de holomorfas es holomorfo cuando el denominador es distinto de $0$.  
Por lo tanto, como $\forall z \neq 0$, 
\[
f(z) = \frac{e^{iz} - 1}{iz} \Rightarrow f(z) \text{ es holomorfa en } \mathbb{C} \setminus \{0\}.
\]

Veamos si $ \displaystyle \lim_{z \to 0} f(z) = 0 $.

\[
\lim_{z \to 0} f(z) = \lim_{z \to 0} \frac{e^{iz} - 1}{iz} = \lim_{z \to 0} \frac{i e^{iz}}{i} = \lim_{z \to 0} e^{iz} = 1.
\]

Por lo tanto, $f$ admite una extensión holomorfa por el teorema de prolongación analítica.  
Esa extensión debe ser continua, por lo tanto:
\[
f(0) = \lim_{z \to 0} f(z) = \lim_{z \to 0} \frac{e^{iz} - 1}{iz} = 1.
\]
Por lo tanto, se cumple efectivamente que $f \in H(\mathbb{C})$.
}
%ejercicio 5
\ejercicio{
    Sea $f$ holomorfa en un dominio $\Omega \subset \mathbb{C}$ tal que $|f(z)| \leq M$ para todo $z \in \Omega$. Demuestra que si $z \in \Omega$, entonces:
    \[
    |f'(z)| \leq \frac{M}{r}
    \]
    donde $r = \text{dist}(z, \text{Fr}(\Omega))$ (distancia a la frontera de $\Omega$).
}{
    Tenemos que $f \in H(\Omega)$ con $\Omega$ dominio abierto y conexo.  
Por lo tanto, sea $z_0 \in \Omega$ y $r = \mathrm{dist}(z_0, \partial \Omega)$.  
Entonces, por el teorema de Taylor se comprueba que
\[
f(z) = \sum_{n=0}^{\infty} c_n (z - z_0)^n,
\]
con 
\[
c_n = \frac{1}{2\pi i} \int_{\partial D(z_0, r')} \frac{f(w)}{(w - z_0)^{n+1}} \, dw,
\quad 0 < r' < \mathrm{dist}(z_0, \partial \Omega).
\]

Por lo tanto:
\[
f'(z_0) = \frac{1}{2\pi i} \int_{\partial D(z_0, r')} \frac{f(w)}{(w - z_0)^2} \, dw.
\]

\[
\Rightarrow \left| f'(z_0)\right| 
\leq \frac{1}{2\pi} \left|\int_{\partial D(z_0, r')} 
\frac{f(w)}{(w - z)^2} dw\right| .
\]

Sea $M = \max_{w \in \partial D(z_0, r')} |f(w)|$. Entonces:
\[
\left| f'(z_0) \right| 
\leq \frac{1}{2\pi} \, \Long(\partial D(z_0, r')) \cdot \frac{M}{(r')^2}
= \frac{2\pi r' M}{2\pi (r')^2} = \frac{M}{r'}.
\]

Con $r' \in (0, r)$, se cumple que:
\[
\left| f'(z_0) \right| \leq \frac{M}{r'} 
\Rightarrow \lim_{r' \to r^-} \left| f'(z_0) \right| \leq lim_{r' \to r^-} \frac{M}{r'i}.
\]
\[
\boxed{\left| f'(z_0) \right| \leq \frac{M}{r}}
\]
\hfill $\square$
}
%ejercicio 6
\ejercicio{
    Sea $f \in H(D(0;1))$ tal que $f(0)=0$, $f'(0)=1$ y $|f(z)| \leq 2$ para todo $z \in D(0;1)$. Demuestra que:
    \[
    |f(z)-z| \leq 1 \quad \forall z \in D(0, 1/2)
    \]
}{
Tenemos $f \in H(D(0,1))$ y $f(0)=0$, $f'(0)=1$. \\[4pt]
Además, $|f(z)| \le 2 \quad \forall z \in D(0,1)$. \\[8pt]

Tenemos que $f \in \mathcal{H}(D(0,1))$, es decir, holomorfa en un abierto, por lo tanto podemos aplicar el tma de Taylor: \\[6pt]

Sea $z_0 = 0 \Rightarrow f(z) = \sum_{n=0}^{\infty} \frac{f^{(n)}(0)}{n!} z^n$, 
$\forall z \in D(0, \bar{r})$. \\[8pt]

Además, 
\[
f(z) = z + \sum_{n=2}^{\infty} \frac{f^{(n)}(0)}{n!} z^n
\]
\[
\Rightarrow |f(z) - z| = \left| \sum_{n=2}^{\infty} \frac{f^{(n)}(0)}{n!} z^n \right|
= \left| \sum_{n=2}^{\infty} C_n z^n \right|
\]
\[
\Rightarrow \sum_{n=2}^{\infty} |C_n| |z|^n
\]
Sabemos que 
\[
C_n = \frac{1}{2\pi i} \int_{\partial D(0,r)} \frac{f(w)}{(w-0)^{n+1}} dw
\]
con $0 < r < 1$. Por lo tanto, consideremos $|z| < r < 1$. \\[8pt]

\[
|C_n| \le \frac{1}{2\pi} |\int_{\partial D(0,r)} \frac{f(w)}{w^{n+1}} dw |
\le \frac{1}{2\pi} \cdot 2\pi r \cdot 
\max_{w \in \partial D(0,r)} \frac{|f(w)|}{r^{n+1}}
= \max_{w \in \partial D(0,r)} \frac{|f(w)|}{r^n}
\]
\[
\le \frac{2}{r^n}
\]
(Estamos considerando $z \in D(0, \tfrac{1}{2})$). \\[10pt]

Entonces:
\[
|f(z) - z| \le \sum_{n=2}^{\infty} |C_n| |z|^n 
\le \sum_{n=2}^{\infty} 2 \left( \frac{1}{2r} \right)^n
\]
\[
= 2 \sum_{n=2}^{\infty} \left( \frac{1}{2r} \right)^n
= \frac{ \left( \frac{1}{2r} \right)^2 }{ 1 - \frac{1}{2r} }
= \frac{ \frac{1}{2} r^3 }{ 2r - 1 }
\quad \xrightarrow[r \to 1^-]{} \quad \frac{1}{2} \le 1 \quad \boxed{\checkmark}
\]

}
%ejercicio 7
\ejercicio{
    Desarrolla en serie de potencias centrada en $z_0$ e indica el radio de convergencia para:
    \begin{enumerate}
        \item $f(z) = \sin z$ en $z_0 = 0$
        \item $f(z) = (1+z)^{-1}$ en $z_0 = i$
    \end{enumerate}
}{
    \textbf{a)} Sabemos que $\sin^2 z \in H(\mathbb{C})$ al ser composición de holomorfas. \\[8pt]

$f \in \mathcal{H}(\mathbb{C}) \Rightarrow f(z) = \displaystyle\sum_{n=0}^{\infty} C_n (z - z_0)^n 
\Rightarrow f(z) = \displaystyle\sum_{n=0}^{\infty} \frac{f^{(n)}(0)}{n!} z^n$ \\[8pt]

con radio de convergencia $R = \infty$. \\[8pt]

\[
\Rightarrow f(z) = \sin^2 z \Rightarrow f(0) = 0
\]
\[
f'(z) = 2\sin z \cos z \Rightarrow f'(0) = 0
\]
\[
f''(z) = 2\cos(2z) \Rightarrow f''(0) = 2
\]
\[
f'''(z) = -4\sin(2z) \Rightarrow f'''(0) = 0
\]
\[
f^{(4)}(z) = -8\cos(2z) \Rightarrow f^{(4)}(0) = -8
\]

Por lo tanto:
\[
f(z) = \sum_{n=1}^{\infty} (-1)^{n+1} \frac{2^{2n-1}}{(2n)!} z^{2n}
\]

\vspace{8pt}
\textbf{b)} Con el mismo procedimiento de antes:

\[
f^{(n)}(z) = \frac{(-1)^n \, n!}{(1+z)^{n+1}}, \quad \forall n \ge 0
\]

\[
\Rightarrow f(z) = \sum_{n=0}^{\infty} \frac{(-1)^n}{(1+i)^{n+1}} (z - i)^n, 
\quad \forall z \in D(i, d(i, \partial\Omega))
\]

con $\Omega = \mathbb{C} \setminus \{-1\}$, 
$\Rightarrow d(i, \partial \Omega) = d(i, -1) = \sqrt{1^2 + 1^2} = \sqrt{2}.$ \\[8pt]

La serie tiene un radio de convergencia $R = \sqrt{2}$ 
(se puede comprobar por el criterio del cociente o raíz).
}
%ejercicio 8
\ejercicio{Definamos un camino $\gamma$ en $\mathbb{C}$ y sea $h: \gamma^* \to \mathbb{C}$ una función continua. Definimos:
    \[
    f(z) := \int_{\gamma} \frac{h(w)}{w-z} \, dw \quad \forall z \in \mathbb{C} \setminus \gamma^*
    \]
    Usa la definición de derivada para probar que $f$ es holomorfa en $\mathbb{C} \setminus \gamma^*$ y que:
    \[
    f'(z) = \int_{\gamma} \frac{h(w)}{(w-z)^2} \, dw
    \]
    \textit{(Indicación: ten en cuenta que $\gamma^*$ es compacto, entonces la distancia está acotada).}}{
    
\textbf{(8)} \\[8pt]
$\gamma$ camino en $\mathbb{C}$ y $h:\gamma^* \to \mathbb{C}$ una función continua. \\[8pt]

Tenemos:
\[
\begin{cases}
f(z) = \displaystyle\int_{\gamma} \frac{h(w)}{w - z} \, dw \quad \forall z \in \mathbb{C} \setminus \gamma^* \\[10pt]
f'(z) = \displaystyle\int_{\gamma} \frac{h(w)}{(w - z)^2} \, dw \, ?
\end{cases}
\]

Tenemos que ver que: sea $z, z_1 \in \mathbb{C} \setminus \gamma^*$ \\[8pt]

\[
\lim_{z \to z_1} \frac{f(z) - f(z_1)}{z - z_1} 
= \int_{\gamma} \frac{h(w)}{(w - z_1)^2} \, dw
\Rightarrow
\]

\[
\lim_{z \to z_1} \frac{f(z) - f(z_1)}{z - z_1}
- \int_{\gamma} \frac{h(w)}{(w - z_1)^2} \, dw = 0
\Rightarrow
\]

\[
\lim_{z \to z_1} 
\int_{\gamma} \left[
\frac{h(w)}{(w - z)(w - z_1)} 
- \frac{h(w)}{(w - z_1)^2}
\right] dw = 0
\]

\[
\int_{\gamma} h(w)
\left(
\frac{1}{(w - z)(w - z_1)} - \frac{1}{(w - z_1)^2}
\right) dw
\]

\[
= |z - z_1| \, \Long(\gamma) \, 
\max_{w \in \gamma^*} 
\left| \frac{h(w)}{(w - z)(w - z_1)} \right|
\quad \text{(*)}
\]

Vemos que $\mathbb{C} \setminus \gamma^*$ es un abierto $\Rightarrow \exists r> 0$ tal que 
$D(z_1, r) \subset \mathbb{C} \setminus \gamma^*$ \\[6pt]
Cogamos entonces $r/2 > 0$. Como $h$ es continua sobre un compacto 
$\gamma^*$ entonces $\exists M > 0$ tal que $|h(w)| \le M \ \forall w \in \gamma^*$ \\[8pt]

\[
\le |z - z_1| \, \Long(\gamma) \, \frac{M}{\left(\frac{r^2}{2}\right)}
\xrightarrow[z \to z_1 ]{ }0\quad \boxed{\checkmark}
\]
}
\ejercicio{}{Queremos ver si $\sum_{n=1}^{\infty} n^{-z}$ es una función holomorfa en el dominio:
\[
A = \{ z \in \mathbb{C} \mid \text{Re}(z) > 1 \}
\]

Sea $f_n(z) = n^{-z}$. Expresamos esto como:
\[
n^{-z} = e^{-z \log n} = e^{-(x+iy)\log n} = e^{-x \log n} e^{-iy \log n} = n^{-x} n^{-iy}
\]
Calculamos el módulo (donde $x = \text{Re}(z)$):
\[
|f_n(z)| = |n^{-x}| \cdot |e^{-iy \log n}| = n^{-x} = \frac{1}{n^x}
\]
Dado que $z \in A$, tenemos $x > 1$.

\subsubsection*{Convergencia uniforme en compactos}
Sea $K$ un compacto de $A$.
Como $K \subset A$ es cerrado y acotado, existe un $\delta > 0$ tal que:
\[
\text{Re}(z) \ge 1 + \delta \quad \forall z \in K
\]
Por lo tanto, acotamos el término general:
\[
|f_n(z)| = \frac{1}{n^{\text{Re}(z)}} \le \frac{1}{n^{1+\delta}}
\]
La serie numérica $\sum_{n=1}^{\infty} \frac{1}{n^{1+\delta}}$ es una serie-p con $p > 1$, por lo que es convergente.

Se cumplen las condiciones del \textbf{Criterio M de Weierstrass}.
Concluimos que $\sum_{n=1}^{\infty} f_n(z)$ converge uniformemente sobre compactos de $A$.

Sea $g_m(z) = \sum_{n=1}^{m} n^{-z}$. Sabemos que $g_m$ es holomorfa (suma finita de holomorfas).
Por el teorema de convergencia de Weierstrass, el límite:
\[
\zeta(z) = \sum_{n=1}^{\infty} n^{-z}
\]
es una función holomorfa en $A$.

La derivada será:
\[
\zeta'(z) = \sum_{n=1}^{\infty} (n^{-z})' = \sum_{n=1}^{\infty} (-\log n) n^{-z}
\]}
